% CSE 719 — Distributed & Cloud Computing
% Lecture 2: Cloud Computing — Model Answers (2020–2024)
% University of Chittagong, 8th Semester B.Sc. Engineering
% Note: 2020 Q1a (cloud+distributed definition) answered in Fundamentals_Solutions.pdf
% Compile: pdflatex Cloud_Solutions.tex (twice for TOC)

\documentclass[12pt, a4paper]{article}

\usepackage[left=2.5cm, right=2.5cm, top=2.8cm, bottom=2.8cm, headheight=15pt]{geometry}
\usepackage{amsmath, amssymb}
\usepackage{booktabs}
\usepackage{array}
\usepackage{xcolor}
\usepackage{enumitem}
\usepackage{fancyhdr}
\usepackage{titlesec}
\usepackage{mdframed}
\usepackage{tabularx}
\usepackage{hyperref}

%----- Colours -----
\definecolor{sectionblue}{RGB}{10,60,110}
\definecolor{boxbg}{RGB}{255,248,240}
\definecolor{answerbg}{RGB}{240,247,255}
\definecolor{notesbg}{RGB}{255,250,230}

%----- Box environments -----
\mdfdefinestyle{qframe}{linecolor=orange!75!black, backgroundcolor=boxbg, roundcorner=4pt, linewidth=1.3pt, innertopmargin=6pt, innerbottommargin=6pt, innerleftmargin=9pt, innerrightmargin=9pt}
\mdfdefinestyle{solframe}{linecolor=sectionblue, backgroundcolor=answerbg, roundcorner=4pt, linewidth=1.3pt, innertopmargin=7pt, innerbottommargin=7pt, innerleftmargin=9pt, innerrightmargin=9pt}
\mdfdefinestyle{notesframe}{linecolor=orange!70!black, backgroundcolor=notesbg, roundcorner=4pt, linewidth=1pt, innertopmargin=5pt, innerbottommargin=5pt, innerleftmargin=8pt, innerrightmargin=8pt}
\newenvironment{qbox}{\begin{mdframed}[style=qframe]}{\end{mdframed}}
\newenvironment{solbox}{\begin{mdframed}[style=solframe]}{\end{mdframed}}
\newenvironment{notesbox}{\begin{mdframed}[style=notesframe]}{\end{mdframed}}

%----- Page style -----
\pagestyle{fancy}
\fancyhf{}
\lhead{\small\color{sectionblue} CSE 719 --- Lecture 2: Cloud Computing}
\rhead{\small\color{sectionblue} Model Solutions}
\cfoot{\small\thepage}
\renewcommand{\headrulewidth}{0.5pt}

%----- Section formatting -----
\titleformat{\section}{\Large\bfseries\color{sectionblue}}{}{0em}{}[\vspace{-2pt}\color{sectionblue}\rule{\linewidth}{0.8pt}]
\titleformat{\subsection}{\large\bfseries\color{sectionblue!85!black}}{}{0em}{}

\newcommand{\qmarks}[1]{\hfill\textbf{\color{orange!70!black}[#1 marks]}}

%=====================================================================
\begin{document}
%=====================================================================
\begin{center}
  {\LARGE\bfseries\color{sectionblue} CSE 719: Distributed \& Cloud Computing}\\[0.5em]
  {\Large\bfseries Lecture 2 --- Cloud Computing}\\[0.3em]
  {\large Model Answers: 2020, 2021, 2022 \& 2024 Papers}\\[0.2em]
  {\small University of Chittagong \quad|\quad 8\textsuperscript{th} Sem B.Sc.\ (Engg.)
  \quad|\quad Sources: Lecture-02 + Unit~3/4 Readings}
\end{center}
\vspace{0.3em}\noindent\rule{\linewidth}{1pt}\vspace{0.5em}
\tableofcontents
\newpage

%=====================================================================
\section{2020 Examination (CSE-813, 52.5 marks)}
%=====================================================================

\textit{Note: 2020 Q1(a) ``Define distributed and cloud systems'' is answered in
\textbf{Fundamentals\_Solutions.pdf} (shared question, see that file).}

%-------
\subsection{Q8(a) --- Define cloud computing; advantages and disadvantages
  \texorpdfstring{\qmarks{2}}{}}
\begin{qbox}Define cloud computing with its advantages and disadvantages.\end{qbox}
\begin{solbox}
\textbf{Definition (NIST):} Cloud computing is a model for enabling \textbf{convenient,
on-demand network access to a shared pool of configurable computing resources} (networks,
servers, storage, applications, services) that can be \textbf{rapidly provisioned and released}
with minimal management effort or service-provider interaction.\\[4pt]

\textbf{Advantages:}
\begin{itemize}[leftmargin=1.4cm, itemsep=1pt]
  \item \textbf{Scalability / Flexibility} --- instantly scale resources up or down; access from
    any internet-connected device.
  \item \textbf{Cost reduction} --- pay-as-you-go; no upfront capital expenditure on hardware.
  \item \textbf{Maintenance offloaded} --- vendor handles patches, updates, monitoring, and
    backups.
  \item \textbf{High availability} --- provider manages fault tolerance; 99.9\%+ SLA uptime
    guaranteed.
\end{itemize}

\textbf{Disadvantages:}
\begin{itemize}[leftmargin=1.4cm, itemsep=1pt]
  \item \textbf{Security \& Privacy} --- sensitive data resides on third-party shared
    infrastructure.
  \item \textbf{Vendor lock-in} --- proprietary APIs make migration to another provider
    costly.
  \item \textbf{Network dependency} --- no internet $\Rightarrow$ no access to applications
    or data.
  \item \textbf{Data migration} --- moving large datasets in/out of a provider is slow and
    expensive.
\end{itemize}
\end{solbox}

%-------
\subsection{Q8(b) --- E-commerce website: cloud services and deployment models
  \texorpdfstring{\qmarks{4}}{}}
\begin{qbox}What cloud services and deployment models would you choose for building an
e-commerce website? Justify your choices.\end{qbox}
\begin{solbox}
\textbf{Cloud Services (by layer):}
\begin{center}\renewcommand{\arraystretch}{1.3}
\begin{tabularx}{\linewidth}{l l X}
\toprule
\textbf{Need} & \textbf{Service Model} & \textbf{Justification} \\
\midrule
Web app hosting & \textbf{PaaS} & Runtime + middleware pre-configured; no OS management; auto-scaling built in (e.g.\ Google App Engine, Heroku). \\
Database / Storage & \textbf{IaaS} & Full control over DB config, replication, and storage volumes (e.g.\ AWS EC2 + EBS). \\
Payment / CRM / Email & \textbf{SaaS} & Off-the-shelf, no custom build --- Stripe (payments), Salesforce (CRM), SendGrid (email). \\
\bottomrule
\end{tabularx}\end{center}

\textbf{Deployment Model:}
\begin{itemize}[leftmargin=1.4cm, itemsep=3pt]
  \item \textbf{Public cloud} (e.g.\ AWS, Azure) for product catalog, storefront, CDN ---
    globally distributed, cost-effective, handles traffic spikes.
  \item \textbf{Private cloud} for payment processing and customer PII ---
    required for regulatory compliance (PCI-DSS, GDPR).
  \item \textbf{Hybrid cloud} overall: routine traffic on public; sensitive data on private;
    \textbf{cloud bursting} to public during sales peaks.
\end{itemize}

\textbf{Justification:} Hybrid maximises cost efficiency while satisfying compliance.
PaaS accelerates development; SaaS avoids reinventing commodity services.
\end{solbox}

%-------
\subsection{Q8(c) --- Cloud aspects for prognostic health management
  \texorpdfstring{\qmarks{2.75}}{}}
\begin{qbox}What aspects of cloud computing are useful for prognostic health management
applications?\end{qbox}
\begin{solbox}
\begin{itemize}[leftmargin=1.4cm, itemsep=3pt]
  \item \textbf{Elastic scalability} --- health sensor data volume is unpredictable; cloud
    elastically provisions compute for ML inference and analytics bursts on demand.
  \item \textbf{Storage elasticity} --- continuous streams of patient/sensor data require
    virtually unlimited storage (e.g.\ Amazon S3, Azure Blob); no pre-purchase.
  \item \textbf{High availability / Reliability} --- provider-managed fault tolerance ensures
    24/7 uptime; critical for real-time health monitoring alerts.
  \item \textbf{Pay-per-use} --- health analytics is batch-heavy; pay only during computation
    phases, not for idle hardware.
  \item \textbf{Security controls} --- encryption at rest/in transit, role-based access control,
    and audit logs required under HIPAA/health regulations.
\end{itemize}
\end{solbox}

%=====================================================================
\section{2021 Examination (CSE 813, 52.5 marks)}
%=====================================================================

%-------
\subsection{Q3(a) --- E-commerce: cloud services and deployment models
  \texorpdfstring{\qmarks{2.75}}{}}
\begin{qbox}Which cloud services and deployment models would you adopt for building an
e-commerce website?\end{qbox}
\begin{solbox}
\textbf{Cloud services:}
\begin{itemize}[leftmargin=1.4cm, itemsep=1pt]
  \item \textbf{IaaS} --- virtual servers for databases; object storage for product images.
  \item \textbf{PaaS} --- web application hosting platform: runtime, middleware, and
    auto-scaling.
  \item \textbf{SaaS} --- payment gateway, email, CRM (no custom build needed).
\end{itemize}
\textbf{Deployment model:}
\begin{itemize}[leftmargin=1.4cm, itemsep=1pt]
  \item \textbf{Hybrid} = Public cloud for storefront and CDN (scalable, low cost) +
    Private cloud for payment/customer PII (regulatory compliance).
  \item Cloud \textbf{bursting}: overflow peak-season traffic to public cloud automatically.
\end{itemize}
\end{solbox}

%-------
\subsection{Q4(c) --- Phases in cloud architecture \texorpdfstring{\qmarks{1}}{}}
\begin{qbox}What are the phases involved in cloud architecture?\end{qbox}
\begin{solbox}
Cloud architecture is built in three successive phases:
\begin{enumerate}[leftmargin=1.6cm, itemsep=2pt]
  \item \textbf{Infrastructure phase (IaaS)} --- physical hardware (servers, storage, network)
    is abstracted via \textbf{virtualization (hypervisor)} into on-demand virtual resources.
  \item \textbf{Platform phase (PaaS)} --- runtime environment, middleware, databases, and
    dev/deployment tools are provisioned on top of the virtual infrastructure.
  \item \textbf{Application phase (SaaS)} --- end-user software is delivered over the internet;
    fully managed by the provider with no user infrastructure responsibility.
\end{enumerate}
\end{solbox}

%-------
\subsection{Q5(a) --- On-demand functionality in cloud
  \texorpdfstring{\qmarks{1.75}}{}}
\begin{qbox}What is on-demand functionality? How is it provided in cloud computing?\end{qbox}
\begin{solbox}
\textbf{On-demand:} Computing resources (compute, storage, bandwidth) are
\textbf{provisioned automatically and immediately} as the user requests them,
\textbf{without any human intervention} from the provider's side.\\[4pt]

\textbf{How it is provided:}
\begin{enumerate}[leftmargin=1.6cm, itemsep=2pt]
  \item Provider maintains \textbf{large shared resource pools} (CPU, RAM, disk, bandwidth).
  \item User requests resources via a \textbf{self-service portal or API} (no phone call to
    provider needed).
  \item The \textbf{hypervisor / orchestration layer} allocates a VM or container within
    seconds.
  \item \textbf{Auto-scaling} monitors load and adds/removes instances dynamically.
  \item \textbf{Metering} tracks consumption; billing reflects actual usage only.
\end{enumerate}
\textit{Key idea:} the user perceives an \textbf{illusion of infinite resources} available
instantly on demand.
\end{solbox}

%-------
\subsection{Q5(b) --- What is a cloud? Large-scale cloud platforms
  \texorpdfstring{\qmarks{1}}{}}
\begin{qbox}What is a cloud? Name some platforms used for large-scale cloud computing.\end{qbox}
\begin{solbox}
\textbf{Cloud (Vaquero et al.):} ``A large pool of easily usable and accessible
\textbf{virtualised resources} (hardware, development platforms, services), dynamically
reconfigured to adjust to a variable load, allowing for optimum resource utilisation via a
\textbf{pay-per-use} model where guarantees are offered by the provider via SLAs.''\\[4pt]

\textbf{Large-scale cloud platforms:}
\begin{itemize}[leftmargin=1.4cm, itemsep=1pt]
  \item \textbf{Amazon EC2 / S3} (AWS) --- dominant public IaaS; virtual machines + object
    storage.
  \item \textbf{Microsoft Azure} --- enterprise-focused public cloud; strong hybrid
    integration.
  \item \textbf{Google Cloud / GCE} --- strong in big data analytics and Kubernetes.
  \item \textbf{EUCALYPTUS} --- open-source \textbf{private} cloud platform, AWS-compatible
    API.
\end{itemize}
\end{solbox}

%-------
\subsection{Q5(c) --- Layers in cloud; cloud vs.\ distributed computing
  \texorpdfstring{\qmarks{3}}{}}
\begin{qbox}Explain the layers in cloud computing and how they work. Differentiate between
cloud and distributed computing.\end{qbox}
\begin{solbox}
\textbf{Layers in cloud (bottom to top):}
\begin{center}\renewcommand{\arraystretch}{1.3}
\begin{tabularx}{\linewidth}{l l X}
\toprule
\textbf{Layer} & \textbf{Model} & \textbf{What it provides} \\
\midrule
1. Physical       & Datacenter    & Servers, switches, power, cooling --- owned by provider. \\
2. Virtualization & Hypervisor    & Abstracts physical hardware into isolated VMs. \\
3. Infrastructure & IaaS          & On-demand VMs, virtual storage, virtual networking. \\
4. Platform       & PaaS          & Runtime, DB, dev tools on top of IaaS. \\
5. Application    & SaaS          & End-user apps delivered via browser; zero local install. \\
\bottomrule
\end{tabularx}\end{center}
\textit{How it works:} User code runs in SaaS; SaaS is built on PaaS (execution environment);
PaaS runs on IaaS (virtual servers); IaaS sits on the physical datacenter.\\[6pt]

\textbf{Cloud vs.\ Distributed Computing:}
\begin{center}\renewcommand{\arraystretch}{1.25}
\begin{tabularx}{\linewidth}{l X X}
\toprule
\textbf{Criterion} & \textbf{Distributed Computing} & \textbf{Cloud Computing} \\
\midrule
Scope        & General model: independent nodes cooperating over a network & Specific implementation delivered as a \textbf{metered utility service} \\
Business model & None inherent & \textbf{Pay-per-use} (OPEX, not CAPEX) \\
Provisioning & Static / pre-allocated by admin & \textbf{Elastic, self-service, on-demand} \\
Management   & User/admin manages nodes & \textbf{Provider} manages all infrastructure \\
Multi-tenancy & Rare & \textbf{Core feature} (shared physical resources) \\
Examples     & DNS, P2P networks, MPI cluster & AWS EC2, Azure, Google Cloud \\
\bottomrule
\end{tabularx}\end{center}
\textbf{Key:} Every cloud is a distributed system, but not every distributed system is a
cloud.
\end{solbox}

%-------
\subsection{Q5(d) --- EUCALYPTUS: what it is, why used, deployment models
  \texorpdfstring{\qmarks{3}}{}}
\begin{qbox}What is EUCALYPTUS? Why is it used? Explain the cloud deployment models.\end{qbox}
\begin{solbox}
\textbf{EUCALYPTUS} (\textit{Elastic Utility Computing Architecture for Linking Your
Programs To Useful Systems}):
\begin{itemize}[leftmargin=1.4cm, itemsep=2pt]
  \item Open-source software framework for building \textbf{private cloud environments} on
    an organisation's existing on-premises hardware.
  \item Implements an \textbf{AWS-compatible API}: tools and scripts written for Amazon EC2
    work on EUCALYPTUS without modification.
  \item Provides IaaS capabilities: VM provisioning, virtual networking, elastic storage ---
    all behind the organisation's firewall.
\end{itemize}

\textbf{Why used:}
\begin{itemize}[leftmargin=1.4cm, itemsep=2pt]
  \item Organisations want cloud elasticity and self-service but \textbf{cannot put data in a
    public cloud} (security policy, data sovereignty regulations, low-latency requirements).
  \item Virtualises in-house servers to reduce idle capacity and hardware sprawl.
  \item AWS-compatible API provides a \textbf{future migration path} to public cloud.
\end{itemize}

\textbf{Cloud Deployment Models:}
\begin{center}\renewcommand{\arraystretch}{1.3}
\begin{tabularx}{\linewidth}{l X}
\toprule
\textbf{Model} & \textbf{Description} \\
\midrule
\textbf{Public}    & Infrastructure owned by a provider (IBM, Google, Amazon); open to paying
                     customers over the internet. \\
\textbf{Private}   & Infrastructure dedicated to a single organisation; may be on-premises
                     (EUCALYPTUS) or hosted exclusively. \\
\textbf{Hybrid}    & Combination of public + private; private for sensitive/regulated workloads,
                     public for elastic bursting. \\
\textbf{Community} & Shared by a consortium of organisations with common concerns (e.g.\ research
                     universities, government agencies). \\
\bottomrule
\end{tabularx}\end{center}
\end{solbox}

%=====================================================================
\section{2022 Examination (CSE 513, 54 marks)}
%=====================================================================

%-------
\subsection{Q1(a) --- Define cloud computing; horizontal vs.\ vertical scaling
  \texorpdfstring{\qmarks{5}}{}}
\begin{qbox}Define cloud computing. Explain horizontal and vertical scaling with suitable
scenarios.\end{qbox}
\begin{solbox}
\textbf{Definition (NIST):} Cloud computing is a model for enabling \textbf{convenient,
on-demand network access to a shared pool of configurable computing resources} that can be
\textbf{rapidly provisioned and released} with minimal management effort or service provider
interaction.\\[6pt]

\textbf{Scaling strategies:}
\begin{center}\renewcommand{\arraystretch}{1.3}
\begin{tabularx}{\linewidth}{l X X}
\toprule
\textbf{Criterion} & \textbf{Vertical Scaling (Scale Up)} & \textbf{Horizontal Scaling (Scale Out)} \\
\midrule
Method           & Add more CPU/RAM/disk to the \textbf{same machine} & Add \textbf{more machines} to the pool \\
Upper limit      & Hardware ceiling of one server & Virtually unlimited (commodity nodes) \\
Downtime         & Often requires restart & None (add/remove nodes live) \\
App changes      & Usually none & May need load balancer, stateless design \\
Fault tolerance  & Single point of failure & No single point of failure \\
Cost             & Expensive (enterprise hardware) & Cheap (commodity servers) \\
\bottomrule
\end{tabularx}\end{center}

\textbf{Scenarios:}
\begin{itemize}[leftmargin=1.4cm, itemsep=3pt]
  \item \textbf{Vertical} --- a legacy monolithic relational database (Oracle) that cannot be
    sharded: give it more RAM/CPU cores. Also suitable as a short-term emergency fix.
  \item \textbf{Horizontal} --- a stateless web/application server tier: add instances behind a
    load balancer. Cloud auto-scaling groups exploit this model for variable traffic.
\end{itemize}
\textbf{Cloud preference:} Horizontal scaling aligns with cloud economics --- use many cheap
VMs; add and remove them elastically.
\end{solbox}

%-------
\subsection{Q1(b) --- E-commerce: cloud services and deployment model
  \texorpdfstring{\qmarks{2}}{}}
\begin{qbox}Which cloud services and deployment model would you suggest for an e-commerce
website?\end{qbox}
\begin{solbox}
\begin{itemize}[leftmargin=1.4cm, itemsep=1pt]
  \item \textbf{IaaS} for database VMs and object storage (product images, order logs).
  \item \textbf{PaaS} for web application tier --- auto-scaled runtime, no OS management.
  \item \textbf{SaaS} for payment gateway, email service, and analytics.
  \item \textbf{Hybrid deployment:} public cloud for storefront and CDN (scalable, cheap);
    private cloud for payment/customer PII (compliance); cloud burst to public at peak sales.
\end{itemize}
\end{solbox}

%-------
\subsection{Q1(c) --- What is multi-tenancy? \texorpdfstring{\qmarks{2}}{}}
\begin{qbox}What is multi-tenancy in cloud computing?\end{qbox}
\begin{solbox}
\textbf{Multi-tenancy:} Multiple customers (\emph{tenants}) share the \textbf{same physical
infrastructure and software instances}, while each tenant's data and configuration remain
\textbf{logically isolated} from every other tenant.\\[4pt]

\textbf{Why it matters:}
\begin{itemize}[leftmargin=1.4cm, itemsep=2pt]
  \item \textbf{Economic model:} provider spreads fixed infrastructure costs across many
    tenants $\Rightarrow$ lower per-customer price.
  \item \textbf{High server utilisation:} shared servers are rarely idle; CPU cycles are sold
    continuously.
\end{itemize}

\textbf{Risks:}
\begin{itemize}[leftmargin=1.4cm, itemsep=1pt]
  \item \textbf{Data leakage} between tenants if VM isolation is misconfigured.
  \item \textbf{Noisy neighbour} effect: one tenant's heavy workload degrades another's
    performance on the same physical host.
\end{itemize}

\textbf{SaaS maturity (Levels 1--4):} Level~1 = separate instance per tenant (weakest, no
sharing). Level~4 = single instance, single database, configurable per-tenant views (most
efficient multi-tenancy).
\end{solbox}

%-------
\subsection{Q2(a) --- EC2 VMs vs.\ physical machines: three benefits to Amazon
  \texorpdfstring{\qmarks{3}}{}}
\begin{qbox}Compare EC2 virtual machines with physical machines. State three benefits Amazon
gains by offering VMs instead of leasing physical machines. What motivated Amazon to start
EC2?\end{qbox}
\begin{solbox}
\textbf{Three benefits to Amazon of VMs over physical machine leasing:}
\begin{enumerate}[leftmargin=1.6cm, itemsep=3pt]
  \item \textbf{Higher server utilisation} --- multiple VMs packed on one physical server
    raise utilisation from $\sim$10--20\% to $>$80\%. More revenue per rack, same power
    and cooling cost.
  \item \textbf{Flexible, fine-grained billing} --- sell fractional capacity (micro, small,
    large, GPU instances) at different price points; maximise yield from each physical
    machine. Physical leasing is coarse-grained (whole server).
  \item \textbf{Strong isolation and safe multi-tenancy} --- hypervisor provides
    hardware-level memory and I/O isolation between tenants; one customer's VM cannot read
    another's memory. Enables renting the same physical server to many customers
    simultaneously.
\end{enumerate}

\textbf{Motivation for EC2:} Amazon sized its datacenters for \textbf{peak holiday shopping
traffic}. During 99\% of the year that capacity was idle. EC2 was launched in 2006 to
\textbf{monetise excess capacity} --- what began as internal infrastructure became the
world's dominant public IaaS platform.
\end{solbox}

%=====================================================================
\section{2024 Examination (CSE-813, 54 marks)}
%=====================================================================

%-------
\subsection{Q5(a) --- Two benefits and two drawbacks of cloud for enterprises
  \texorpdfstring{\qmarks{2.5}}{}}
\begin{qbox}State two benefits and two drawbacks of cloud computing for enterprises.\end{qbox}
\begin{solbox}
\textbf{Benefits:}
\begin{enumerate}[leftmargin=1.6cm, itemsep=3pt]
  \item \textbf{Reduced capital expenditure} --- pay-as-you-go converts server buying (CAPEX)
    into a monthly utility bill (OPEX); small and medium enterprises access the infrastructure
    of corporate giants with no upfront investment.
  \item \textbf{Elastic scalability on demand} --- instantly provision extra compute for
    demand spikes (product launches, end-of-quarter runs) without buying hardware that sits
    idle the rest of the year.
\end{enumerate}

\textbf{Drawbacks:}
\begin{enumerate}[leftmargin=1.6cm, itemsep=3pt]
  \item \textbf{Security and privacy risk} --- sensitive corporate data (IP, financials,
    customer records) resides on third-party shared infrastructure; data breaches or
    regulatory non-compliance (GDPR) expose the enterprise to liability.
  \item \textbf{Vendor lock-in} --- heavy use of provider-specific services (AWS Lambda,
    Azure Cosmos DB, Google BigQuery) makes migration to a competitor slow and expensive;
    the enterprise loses negotiating leverage over pricing and SLAs.
\end{enumerate}
\end{solbox}

%-------
\subsection{Q5(b) --- Virtualization in cloud; two benefits to providers
  \texorpdfstring{\qmarks{2.5}}{}}
\begin{qbox}Explain the fundamental concept of virtualization in cloud computing. State two
benefits of virtualization to cloud providers for resource management.\end{qbox}
\begin{solbox}
\textbf{Virtualization in cloud:}
\begin{itemize}[leftmargin=1.4cm, itemsep=2pt]
  \item A \textbf{hypervisor} (Virtual Machine Monitor) runs directly on physical hardware
    and \textbf{abstracts hardware resources} into multiple independent \textbf{Virtual
    Machines (VMs)}.
  \item Each VM has its own OS, CPU allocation, memory, and disk --- fully \textbf{isolated}
    from other VMs on the same physical server.
  \item The stack: \textbf{Physical hardware}
    $\rightarrow$ \textbf{Hypervisor}
    $\rightarrow$ \textbf{VMs (Guest OS + Applications)}.
  \item Types: \textit{Full virtualisation} (Type-1 bare-metal: VMware ESXi, Xen, KVM) and
    \textit{para-virtualisation} (guest OS is aware of the hypervisor, lower overhead).
\end{itemize}

\textbf{Two benefits to cloud providers:}
\begin{enumerate}[leftmargin=1.6cm, itemsep=3pt]
  \item \textbf{Dramatically higher server utilisation} --- packing multiple VMs on one
    physical server raises utilisation from $\sim$20\% to $>$80\%, reducing the number of
    physical machines (and their power/cooling cost) for a given workload.
  \item \textbf{Rapid, flexible provisioning} --- a new VM is created in seconds from a
    template; the provider scales customer capacity instantly without shipping or installing
    new hardware, enabling true on-demand elasticity.
\end{enumerate}
\end{solbox}

%-------
\subsection{Q6(b) --- IaaS vs.\ PaaS trade-offs for e-commerce migration
  \texorpdfstring{\qmarks{3}}{}}
\begin{qbox}An e-commerce company is migrating to the cloud. Compare IaaS and PaaS on:
development speed, operational overhead, vendor lock-in, and scalability.\end{qbox}
\begin{solbox}
\begin{center}\renewcommand{\arraystretch}{1.35}
\begin{tabularx}{\linewidth}{l X X}
\toprule
\textbf{Criterion} & \textbf{IaaS} & \textbf{PaaS} \\
\midrule
\textbf{Dev speed}  & Slow --- team must install and configure OS, web server, runtime,
                      middleware & Fast --- platform pre-configured; developer deploys code directly \\[4pt]
\textbf{Op.\ overhead} & High --- team manages OS patches, security hardening, scaling
                      configuration & Low --- provider manages OS, middleware, runtime updates \\[4pt]
\textbf{Vendor lock-in} & Low --- VMs are portable; standard OS (Linux/Windows) is the only
                      dependency & High --- proprietary runtimes, build APIs (e.g.\ AWS Elastic
                      Beanstalk, Heroku) \\[4pt]
\textbf{Scalability} & Manual --- team configures auto-scaling rules and load balancers
                      explicitly & Automatic --- built-in horizontal scaling with no extra
                      configuration \\[4pt]
\midrule
\textbf{Best for} & Database servers, legacy apps, custom network topologies & Stateless web/app
                    tier, rapid prototyping, microservices \\
\bottomrule
\end{tabularx}\end{center}

\textbf{Recommendation for e-commerce:} Use \textbf{PaaS for the application tier} (fast
deployment, auto-scaling handles traffic spikes) and \textbf{IaaS for the database layer}
(full control over DB configuration and storage volumes; lower lock-in risk for critical
data).
\end{solbox}

%-------
\subsection{Q6(c) --- IaaS cost calculation: reserved vs.\ on-demand instances
  \texorpdfstring{\qmarks{3}}{}}
\begin{qbox}Explain reserved instances vs.\ on-demand instances in IaaS pricing. Calculate
the cost for a predictable base load and a variable peak workload.\end{qbox}
\begin{solbox}
\textbf{Reserved instances:} Pay an upfront commitment (1-year or 3-year term) in exchange
for a \textbf{discounted hourly rate} (typically 30--70\% cheaper than on-demand). Best for
\textbf{predictable, always-on} workloads.\\[4pt]

\textbf{On-demand instances:} Pay the full hourly rate with \textbf{no commitment}; terminate
any time. Best for \textbf{variable, unpredictable} workloads, testing, or seasonal
peaks.\\[6pt]

\textbf{Cost calculation example (monthly):}
\begin{center}\renewcommand{\arraystretch}{1.3}
\begin{tabularx}{0.92\linewidth}{l X r}
\toprule
\textbf{Workload} & \textbf{Detail} & \textbf{Monthly cost} \\
\midrule
Base load (reserved)  & 10 instances $\times$ 720 hr $\times$ \$0.05/hr & \$360 \\
Peak load (on-demand) & 5 extra instances $\times$ 100 hr $\times$ \$0.10/hr & \$50 \\
\midrule
\textbf{Total} & & \textbf{\$410} \\
\bottomrule
\end{tabularx}\end{center}

\textbf{Decision rule:}
\begin{itemize}[leftmargin=1.4cm, itemsep=1pt]
  \item Instance runs $>$~50\% of the time $\Rightarrow$ \textbf{reserved is cheaper}.
  \item Instance runs $<$~50\% of the time $\Rightarrow$ \textbf{on-demand is cheaper}.
  \item Optimal strategy: \textbf{reserve the predictable base load}; use on-demand for
    traffic spikes.
\end{itemize}
\end{solbox}

%=====================================================================
\section{Quick Summary --- Exam Patterns \& Must-Know}
%=====================================================================
\begin{notesbox}
\begin{itemize}[leftmargin=1.3cm, itemsep=3pt]
  \item \textbf{Two definitions to memorise:} NIST (``on-demand, shared pool, rapidly
    provisioned'') + Vaquero (``large pool of virtualised resources, pay-per-use, SLA'').
  \item \textbf{E-commerce deployment (2020 Q8b, 2021 Q3a, 2022 Q1b) --- verbatim 3 years.}
    Lock in: IaaS (DB/storage) + PaaS (app tier) + SaaS (payment/email); Hybrid model (public
    storefront + private PII data). Prepare a 2-minute mental rehearsal.
  \item \textbf{Pros/cons} (2020 Q8a, 2024 Q5a): Pros = Flexibility, Scalability, Cost, Maintenance
    offloaded, Utilisation, Availability. Cons = Security, Privacy, Vendor lock-in, Network
    dependency, Migration.
  \item \textbf{Horizontal vs vertical scaling (2022 Q1a):} Vertical = one big machine; Horizontal
    = many small machines. Cloud favours horizontal. Know the 6-row comparison table.
  \item \textbf{Multi-tenancy (2022 Q1c):} Shared physical infra + logical isolation. Risks:
    data leakage, noisy neighbour.
  \item \textbf{Virtualization (2024 Q5b):} Physical $\rightarrow$ Hypervisor $\rightarrow$
    VMs. Two provider benefits: utilisation + rapid provisioning.
  \item \textbf{IaaS vs PaaS (2024 Q6b):} Dev speed / overhead / lock-in / scalability ---
    know the 4-row comparison table cold.
  \item \textbf{EUCALYPTUS (2021 Q5d):} Open-source private cloud, AWS-compatible API. Used
    when data cannot leave on-premises for security/compliance reasons.
  \item \textbf{Cloud vs Distributed (2021 Q5c):} Cloud IS a distributed system but adds
    pay-per-use, elasticity, multi-tenancy, provider-managed infrastructure.
  \item \textbf{EC2 motivation (2022 Q2a):} Amazon had excess capacity from holiday peaks ---
    monetised it. Three VM benefits: utilisation + fine-grained billing + isolation.
\end{itemize}
\end{notesbox}

\end{document}
